\documentclass[11pt]{article}

\usepackage[utf8]{inputenc}
\usepackage[T1]{fontenc}
\usepackage{lmodern}
\usepackage[margin=2.7cm]{geometry}
\usepackage{amsmath,amssymb,amsthm,mathtools}
\usepackage{booktabs}
\usepackage[hidelinks]{hyperref}
\usepackage[capitalise,noabbrev]{cleveref}
\usepackage{microtype}
\usepackage{url}
\usepackage[shortlabels]{enumitem}

\newtheorem{theorem}{Theorem}
\newtheorem{lemma}[theorem]{Lemma}
\newtheorem{proposition}[theorem]{Proposition}
\newtheorem{corollary}[theorem]{Corollary}
\theoremstyle{definition}
\newtheorem{definition}[theorem]{Definition}
\theoremstyle{remark}
\newtheorem{remark}[theorem]{Remark}

\newcommand{\Lq}{\mathsf L}
\newcommand{\Gq}{\mathsf G}
\newcommand{\Aq}{\mathsf A}
\newcommand{\Bq}{\mathsf B}
\newcommand{\Fq}{\mathsf F}

\title{Early rows of the reflected triangles, firing classes,\\
and the five-state near misses}
\author{Research notes accompanying the note\\
\emph{Two necessary conditions on minimal-time solutions}\\
\emph{of the firing squad problem}}
\date{October 2026}

\begin{document}
\maketitle

\noindent Notation and numbering of results refer to the note: the half-line
$C_\infty$, the line $C_n$, the reflected triangle $R_n$ in the relative
coordinates $\tau=t-(n-1)$, $\kappa=n-i$, so that
$R_n=\{(\tau,\kappa):0\le\kappa\le\tau\le n-1\}$, the input word
\[
J_n(\kappa)=\bigl(\Phi_{n-3+\kappa}(2\kappa-2),\;\Phi_{n-2+\kappa}(2\kappa-1)\bigr),
\qquad 0\le\kappa\le n-1,
\]
the depth rows $\Phi_t(j)=C_\infty(t,t+1-j)$ with the recurrence
\begin{equation}\label{eq:comoving}
\Phi_{t+1}(j)=\delta\bigl(\Phi_t(j),\Phi_t(j-1),\Phi_t(j-2)\bigr)\qquad
(0\le j\le t+1),
\end{equation}
where $\Phi_t(-1)=\Phi_t(-2)=\Lq$ and $\Phi_t(t+1)=\ast$ (equation~(3) of the
note), and the partial map $\Psi$ of the determinacy lemma (Lemma~3.1 of the
note). We write $R_n(\tau,\kappa):=C_n(n-1+\tau,n-\kappa)$. All statements are
proved below; the computer-assisted parts are simulations of explicitly given
rule tables and use no SAT solver. They are reproduced by the programs
\texttt{asymp.py} and \texttt{failall.py} in the directory
\texttt{fssp5/research} of the repository.

\section{Early rows}\label{sec:early}

For $T\ge0$ let $E(T):=\{(\tau,\kappa):0\le\kappa\le\tau\le T\}$. If $n\ge T+2$,
then $E(T)\subseteq R_n$ and no cell of $E(T)$ lies in the firing row
$\tau=n-1$ of $R_n$.

\begin{lemma}[early cells]\label{lem:early}
Let $\delta$ be a rule, $T\ge0$, and $n,n'\ge T+2$ with $J_n(\kappa)=J_{n'}(\kappa)$
for $0\le\kappa\le T+1$. Then for every $(\tau,\kappa)\in E(T)$,
$\Psi(J_n,n;\tau,\kappa)$ is defined if and only if
$\Psi(J_{n'},n';\tau,\kappa)$ is defined, and then the two values are equal. In
particular, the line of length $n$ has a cell in state $\Fq$ at a relative time
$\tau\le T$ if and only if the line of length $n'$ has one at the same relative
position.
\end{lemma}

\begin{proof}
By the definition of $\Psi$ in the proof of Lemma~3.1 of the note,
\[
\Psi(J,m;\tau,\kappa)=\delta(x_{\kappa+1},x_\kappa,x_{\kappa-1}),
\]
where $x_{\kappa'}$ is the border $\ast$ if $\kappa'\in\{-1,m\}$, a component of
$J(\kappa')$ if $\tau-1\in\{\kappa'-2,\kappa'-1\}$, and
$\Psi(J,m;\tau-1,\kappa')$ otherwise (then $0\le\kappa'\le\tau-1$). Unfolding
this recursion from $(\tau,\kappa)\in E(T)$ visits only pairs
$(\tau'',\kappa'')$ with $0\le\kappa''\le\tau''\le T$ and reads the symbols
$x_{\kappa'}$ with $\kappa'\le\kappa''+1\le T+1$. The case $\kappa'=m$ would
require $m\le T+1$, which is excluded for $m\in\{n,n'\}$; hence the left border
is never read, and the input word is read only at positions $\kappa'\le T+1$,
where $J_n$ and $J_{n'}$ agree. The two recursions are therefore identical step
by step, including the points where an argument equal to $\Fq$ makes $\delta$
undefined. The last statement follows because the cells of the line of length
$n$ outside $R_n$ coincide with cells of $C_\infty$ (Lemma~2.2 of the note) and
are never $\Fq$ before the line has fired, while the cells of $R_n$ at the
relative times $\tau\le T$ are exactly the cells of $E(T)$.
\end{proof}

\begin{lemma}[periodicity certificate]\label{lem:certificate}
Let $\delta$ be a rule whose half-line does not fire, $D\ge0$, and
$V_t:=(\Phi_t(0),\dots,\Phi_t(D))$ for $t\ge D$. There is a map $G_D$, depending
on $\delta$ and $D$ only, with $V_{t+1}=G_D(V_t)$ for all $t\ge D$. Consequently,
if $V_{t_0}=V_{t_0+p}$ for some $t_0\ge D$ and $p\ge1$, then $V_t=V_{t+p}$ for
all $t\ge t_0$.
\end{lemma}

\begin{proof}
For $t\ge D$ and $0\le j\le D$ we have $j\le t$, so the three arguments in
\eqref{eq:comoving} are $\Phi_t(j)$, $\Phi_t(j-1)$, $\Phi_t(j-2)$ with indices at
most $D$ (and $\Lq$ for negative indices); none of them is the border symbol
$\Phi_t(t+1)=\ast$. Hence $G_D(v)_j:=\delta(v_j,v_{j-1},v_{j-2})$, with
$v_{-1}=v_{-2}=\Lq$, does it; it is well defined on the vectors that occur
because the half-line does not fire. The second statement follows by induction
on $t$.
\end{proof}

\begin{corollary}[finitely many early-row problems]\label{cor:reduce}
Let $\delta$, $T\ge0$, $D:=2T+1$, $t_0\ge D$ and $p\ge1$ satisfy
$V_{t_0}=V_{t_0+p}$. Then $J_n(\kappa)=J_{n+p}(\kappa)$ for all
$0\le\kappa\le T+1$ and all $n\ge t_0+3$. Hence, for every $N_0\ge T+2$, each
line of length $n\ge N_0$ has the same cells at the relative times $\tau\le T$
(in the sense of \cref{lem:early}) as one of the lines of the finitely many
lengths $N_0\le n'<\max(N_0,t_0+3)+p$.
\end{corollary}

\begin{proof}
For $0\le\kappa\le T+1$ the two components of $J_n(\kappa)$ are depth-row cells
of depths $2\kappa-2,2\kappa-1\le 2T+1=D$ at the times $n-3+\kappa$ and
$n-2+\kappa$, which are at least $n-3\ge t_0$ (for $\kappa=0$ the lower
component is $\Lq$ by convention, and the upper one is $C_\infty(n-2,n)=\Lq$ by
Lemma~2.1 of the note). By \cref{lem:certificate} these cells are unchanged when
$n$ is replaced by $n+p$. If $n\ge\max(N_0,t_0+3)+p$, subtracting a suitable
multiple of $p$ gives $n'$ in $[\max(N_0,t_0+3),\max(N_0,t_0+3)+p)$ with
$J_{n'}=J_n$ on $[0,T+1]$, and \cref{lem:early} applies since $n,n'\ge N_0\ge
T+2$.
\end{proof}

\begin{remark}\label{rem:limit}
For a minimal-time solution, every line of length $n\ge T+2$ has no $\Fq$ at
the relative times $\tau\le T$. If the depth column $V_t$ for $D=2T+1$ is
periodic from $t_0$ with period $p$, the early rows of $R_n$ for $n\ge t_0+3$
depend only on $n\bmod p$: they are the first $T+1$ rows of finitely many
\emph{limit triangles}, which grow from the right border on the periodic
background formed by the depth rows near the front. A minimal-time solution
needs every limit triangle to stay free of $\Fq$ in all rows, just as the
half-line itself never fires. The near misses of the note violate exactly this,
as the next proposition shows.
\end{remark}

\begin{proposition}\label{prop:nearmiss}
None of the rules $\delta_{14}$, $\delta_{14}^{\mathrm b}$,
$\delta_{14}^{\mathrm c}$, $\delta_{14}^{\mathrm d}$, $\delta_{14}^{\mathrm e}$
(files \texttt{results/delta14*.txt} of the repository) synchronizes a line of
length $n\ge15$ at time $2n-2$. Each of them therefore synchronizes exactly the
lines of lengths $2,\dots,14$ in minimal time.
\end{proposition}

\begin{proof}[Proof (computer-assisted, by simulation)]
For each rule, \cref{tab:certificates} gives $T$ and a repetition
$V_{t_0}=V_{t_0+p}$ of the depth column for $D=2T+4\ge2T+1$, found by
simulating the half-line. Each rule's half-line uses only the transitions of its
germ and does not fire on the computed cells, so \cref{lem:certificate} applies.
By direct simulation:
\begin{enumerate}[(i)]
\item every line of length $15\le n<T+3$ fails to synchronize at time $2n-2$
(early or incomplete firing);
\item for every $n$ with $\max(15,T+3)\le n<t_0+3+p$, the line of length $n$, or
a line of a length $n'$ in the same range with $J_{n'}=J_n$ on $[0,T+1]$, has a
cell in state $\Fq$ at a relative time $\tau\le T$.
\end{enumerate}
By \cref{lem:early} and \cref{cor:reduce} (with $N_0=\max(15,T+3)$), every line
of length $n\ge\max(15,T+3)$ has an $\Fq$ at a relative time $\tau\le T<n-1$,
that is, before time $2n-2$.
The lines of lengths $2,\dots,14$ are synchronized by these rules
(Proposition~5.2 of the note and its text).
\end{proof}

\begin{table}[ht]
\centering\small
\begin{tabular}{@{}lccccc@{}}
\toprule
rule & $T$ & $t_0$ & $p$ & lengths simulated & first $\Fq$ at relative time\\
\midrule
$\delta_{14}$ & $10$ & $51$ & $8$ & $35$ $(15,\dots,49)$ & $7,\dots,10$\\
$\delta_{14}^{\mathrm b}$ & $14$ & $66$ & $3$ & $39$ $(15,\dots,53)$ & $8,\dots,14$\\
$\delta_{14}^{\mathrm c}$ & $22$ & $257$ & $192$ & $413$ $(15,\dots,427)$ & $7,\dots,22$\\
$\delta_{14}^{\mathrm d}$ & $12$ & $95$ & $48$ & $116$ $(15,\dots,130)$ & $6,\dots,12$\\
$\delta_{14}^{\mathrm e}$ & $14$ & $56$ & $2$ & $28$ $(15,\dots,42)$ & $9,\dots,14$\\
\bottomrule
\end{tabular}
\caption{Certificates for \cref{prop:nearmiss}: the depth column of
depths $0,\dots,2T+4$ repeats with period $p$ from time $t_0$; one length per
distinct signature (input cells at $\kappa\le T+2$ and the return-chain values
of $C_\infty$, a slightly finer partition than needed) is simulated, together
with every length $15\le n<T+3$. Program: \texttt{failall.py RULE T}.}
\label{tab:certificates}
\end{table}

\paragraph{The early-rows test} Let $g$ be a germ that determines $C_\infty$
on the cells used below, and let the depth column of $C_\infty$ for $D=2T+1$
satisfy $V_{t_0}=V_{t_0+p}$. If a rule $\delta\supseteq g$ is a minimal-time
solution, then for each of the finitely many lengths $n$ of
\cref{cor:reduce} no cell of $R_n$ with $\tau\le\min(T,n-2)$ is $\Fq$, and
$R_n(\tau,\tau)\ne C_\infty(n-1+\tau,n-\tau)$ for $\tau\le\min(T,n-1)$ by
Lemma~2.3 of the note. By \cref{cor:reduce}, these conditions for the listed
lengths imply the same conditions for every $n\ge N_0$: no further length adds
a constraint on the early rows. Combined with the complete triangles of the
lengths $2,\dots,N_0-1$, they form a finite formula; if it has no model
extending $g$, no minimal-time solution contains $g$. This test is the encoding
\texttt{early:n:T} of \texttt{gfix.py}, with the lengths produced by
\texttt{asymp.py}.

\section{Uniform pre-firing configurations}\label{sec:classes}

\begin{lemma}\label{lem:class}
Let $X\in W$, $N\ge3$, and let $\delta$ be a minimal-time solution such that for
every $3\le n\le N$ the configuration of the line of length $n$ at time $2n-3$
is $X^n$. Define $\delta'$ by $\delta'(x,y,z):=\Lq$ for every $(x,y,z)\in
W^3\setminus\{(X,X,X)\}$ with $\delta(x,y,z)=\Fq$, and $\delta':=\delta$
elsewhere. Then $\delta'(X,X,X)=\Fq$, no other neighbourhood without a border is
mapped to $\Fq$ by $\delta'$, and $\delta'$ coincides with $\delta$ on every
neighbourhood that occurs in the half-line or in a line of length $n\le N$
before time $2n-2$. In particular, $\delta'$ satisfies every constraint that
the lines of lengths $2,\dots,N$ and the half-line impose on $\delta$.
\end{lemma}

\begin{proof}
The rule $\delta'$ satisfies the quiescence condition~(1) of the note, since
$\delta(\Lq,\Lq,\Lq)=\Lq$. The line of length $3$ fires from $XXX$, so
$\delta(X,X,X)=\Fq=\delta'(X,X,X)$. Let $(x,y,z)\in W^3\setminus\{(X,X,X)\}$
with $\delta(x,y,z)=\Fq$. It does not occur in the half-line, which never fires
(equation~(2) of the note). If it occurred as the neighbourhood of a cell of the
line of length $n\le N$ at a time $t<2n-2$, that cell would be $\Fq$ at time
$t+1$; this is impossible for $t+1<2n-2$, and for $t=2n-3$ the configuration is
$X^n$, whose neighbourhoods without a border are all $(X,X,X)$ (the line of
length $2$ has none). Changing the values of neighbourhoods that never occur
changes neither the half-line nor these lines.
\end{proof}

So if the constraints of the lengths $2,\dots,N$ (with the half-line encoded to
any anti-diagonal), together with ``the only neighbourhood without a border
that gives $\Fq$ is $(X,X,X)$'', are unsatisfiable, then no minimal-time
solution fires from $X^n$ at every length $3\le n\le N$. This is option
\texttt{--Pint XXX} of \texttt{gcegar.py}; option \texttt{--upf X} adds the
implied unit clauses ``every cell at time $2n-3$ is $X$''. The state
$X=\Lq$ is impossible, because $\delta(\Lq,\Lq,\Lq)=\Lq$, and $X=\Bq$ is the
case $X=\Aq$ after exchanging the auxiliary states. Mazoyer's solution fires
from $\Gq^n$, and so do the rules satisfying Balzer's condition (B3) in the
strong reading.

\section{A pumping lemma for the right band}\label{sec:column}

The left-border lemma (Lemma~4.1 of the note) runs a band of chains along an
anti-diagonal. The same pigeonhole argument applies to the cells next to the
right border, read in the direction of time. It gives a necessary condition on
the reflected triangle itself.

\begin{lemma}[right band]\label{lem:column}
Let $\delta$ have $k$ states and synchronize the line of length $n$ at time
$2n-2$. Let $1\le d\le\tau_0\le n-2$ and $\ell\ge1$, and suppose that the
column $\tau\mapsto R_n(\tau,d)$ is $\ell$-periodic on $[\tau_0,n-2]$. Then
$n-1-\tau_0\le(k-1)^d\ell$.
\end{lemma}

\begin{proof}
For $d-1\le\tau\le n-1$ put $\mathbf b(\tau):=(R_n(\tau,0),\dots,R_n(\tau,d-1))$.
Since $R_n(\tau+1,\kappa)=\delta(R_n(\tau,\kappa+1),R_n(\tau,\kappa),
R_n(\tau,\kappa-1))$ with $R_n(\tau,-1)=\ast$, there is a map $\Gamma$, depending
on $\delta$ and $d$ only, with
$\mathbf b(\tau+1)=\Gamma(\mathbf b(\tau),R_n(\tau,d))$ for $d\le\tau\le n-2$; its
arguments lie in $W^d\times W$ because no cell fires before time $2n-2$. At
$\tau=n-1$ the band is $\Fq^d$. Suppose $n-1-\tau_0>(k-1)^d\ell$. The pairs
$(\mathbf b(\tau),\tau\bmod\ell)$, $\tau_0\le\tau\le n-2$, take at most
$(k-1)^d\ell$ values, so there are $\tau_0\le\tau_1<\tau_2\le n-2$ with
$\mathbf b(\tau_1)=\mathbf b(\tau_2)$ and $\ell\mid h:=\tau_2-\tau_1$. For
$0\le s\le n-2-\tau_2$ the driving cells $R_n(\tau_1+s,d)$ and $R_n(\tau_2+s,d)$
lie in $[\tau_0,n-2]$ and differ by a multiple of the period, so they are
equal, and induction on $s$ gives
$\mathbf b(\tau_1+s)=\mathbf b(\tau_2+s)$ for $0\le s\le n-1-\tau_2$. For
$s=n-1-\tau_2$ this reads $\mathbf b(n-1-h)=\mathbf b(n-1)=\Fq^d$, although
$n-1-h<n-1$: a contradiction.
\end{proof}

Unlike the two barriers of the note, the hypothesis concerns cells of $R_n$, not
of the half-line, so the lemma is no test on germs. The partial-triangle tests
(cones of firing cells, \texttt{gfix.py}) imply it for every length they
encode.

\section{The complexity-15 half-line}\label{sec:labellings}

The two germs of complexity $15$ that survive the tests of the note are those of
$\delta_{14}^{\mathrm c}$ and $\delta_{14}^{\mathrm d}$. Let $\sigma$ be the
renaming $\Aq\mapsto\Gq$, $\Bq\mapsto\Aq$, $\Gq\mapsto\Bq$ of the states of
$\delta_{14}^{\mathrm d}$ ($\Lq$ and $\Fq$ fixed).

\begin{lemma}\label{lem:labellings}
$\sigma$ maps the thirteen transitions of the germ of $\delta_{14}^{\mathrm d}$
that occur at times $t\ge1$ onto the thirteen transitions of the germ of
$\delta_{14}^{\mathrm c}$ that occur at times $t\ge1$. Let $n\ge3$. In the names
of $\delta_{14}^{\mathrm d}$, the line of length $n$ is $\Aq\,\Bq\,\Lq^{n-2}$ at
time $1$ for both germs. Hence the constraints that the line of length $n$
imposes on a completion of the germ of $\delta_{14}^{\mathrm d}$ and,
transported by $\sigma^{-1}$, on a completion of the germ of
$\delta_{14}^{\mathrm c}$ are the same. The only exceptions are the four entries
$(\ast,\Gq,\Lq)$, $(\Gq,\Lq,\Lq)$, which the germ of $\delta_{14}^{\mathrm d}$
fixes to $\Aq,\Bq$, and $(\ast,\Aq,\Lq)$, $(\Aq,\Lq,\Lq)$, which the transported
germ of $\delta_{14}^{\mathrm c}$ fixes to $\Aq,\Bq$.
\end{lemma}

\begin{proof}
The germ of $\delta_{14}^{\mathrm d}$ is $\ast\Gq\Lq{>}\Aq$, $\Gq\Lq\Lq{>}\Bq$,
$\ast\Aq\Bq{>}\Aq$, $\Aq\Bq\Lq{>}\Bq$, $\Bq\Lq\Lq{>}\Bq$, $\Aq\Bq\Bq{>}\Bq$,
$\Bq\Bq\Lq{>}\Lq$, $\Bq\Lq\Bq{>}\Bq$, $\Lq\Bq\Lq{>}\Aq$, $\Bq\Bq\Aq{>}\Lq$,
$\Bq\Aq\Bq{>}\Bq$, $\Lq\Bq\Bq{>}\Lq$, $\Bq\Bq\Bq{>}\Bq$, $\Lq\Bq\Aq{>}\Bq$,
$\Aq\Bq\Aq{>}\Aq$. That of $\delta_{14}^{\mathrm c}$ is $\ast\Gq\Lq{>}\Gq$,
$\Gq\Lq\Lq{>}\Aq$, $\ast\Gq\Aq{>}\Gq$, $\Gq\Aq\Lq{>}\Aq$, $\Aq\Lq\Lq{>}\Aq$,
$\Gq\Aq\Aq{>}\Aq$, $\Aq\Aq\Lq{>}\Lq$, $\Aq\Lq\Aq{>}\Aq$, $\Lq\Aq\Lq{>}\Gq$,
$\Aq\Aq\Gq{>}\Lq$, $\Aq\Gq\Aq{>}\Aq$, $\Lq\Aq\Aq{>}\Lq$, $\Aq\Aq\Aq{>}\Aq$,
$\Lq\Aq\Gq{>}\Aq$, $\Gq\Aq\Gq{>}\Gq$. Applying $\sigma$ to the last thirteen
transitions of the first list gives the last thirteen of the second. In both
germs the first two transitions are used only at $t=0$ (the general is the only
cell in state $\Gq$ of $\delta_{14}^{\mathrm d}$'s half-line, and
$(\ast,\Gq,\Lq)$, $(\Gq,\Lq,\Lq)$ do not occur in $\delta_{14}^{\mathrm c}$'s half-line
at $t\ge1$, by the correspondence). At time $1$ the line of length $n\ge3$ is
$(\delta(\ast,\Gq,\Lq),\delta(\Gq,\Lq,\Lq),\Lq,\dots,\Lq)$, which is
$\Aq\,\Bq\,\Lq^{n-2}$ for $\delta_{14}^{\mathrm d}$ and $\Gq\,\Aq\,\Lq^{n-2}$, that
is $\sigma(\Aq\,\Bq\,\Lq^{n-2})$, for $\delta_{14}^{\mathrm c}$. From time $1$
on, the line evolves by the table alone, so a completion $\delta$ of the first
germ and the completion $\sigma\circ\delta\circ\sigma^{-1}$ of the second (after
the four entries above are set as stated) produce corresponding lines of every
length $n\ge3$. The line of length $2$ reads $(\Gq,\Lq,\ast)$ at time $0$ in
both germs, which corresponds to two different entries.
\end{proof}

Consequently the completion problems of the two labellings for the lengths
$3,\dots,N$, with the four entries left free, have a common relaxation; if it
is unsatisfiable, both germs are refuted and every five-state minimal-time
solution has $c_{78}\ge16$.

\end{document}
